\section{Arhitectura operațională AES-GCM}
\label{anexa:cripto_aes-gcm}
Pentru criptarea autentificată protocolul folosește \textit{AES} în modul \textit{GCM}. Procesul este împărțit în fluxul de criptare/decriptare, pe baza funcției GCTR, și calcularea etichetei de autentificare prin funcția GHASH. Operațiile acestea sunt prezentate vizual în Figurile~\ref{fig:aes_enc}--\ref{fig:aes_ghash}, urmate de definițiile matematice, informațiile fiind preluate din \textit{NIST Special Publication 800-38D} ~\cite{nist_aesgcm}.

\begin{figure}[htbp]
    \centering
    \begin{minipage}[t]{0.48\textwidth}
        \centering
        \includegraphics[width=\textwidth]{aes-gcm_encryption.png}
        \caption{Criptare AES-GCM}
        \label{fig:aes_enc}
    \end{minipage}\hfill
    \begin{minipage}[t]{0.48\textwidth}
        \centering
        \includegraphics[width=\textwidth]{aes-gcm_decryption.png}
        \caption{Decriptare AES-GCM}
        \label{fig:aes_dec}
    \end{minipage}
\end{figure}

\vspace{0.2cm}

\begin{figure}[htbp]
    \centering
    \begin{minipage}[t]{0.48\textwidth}
        \centering
        \includegraphics[width=\textwidth]{gctr_function.png}
        \caption{Operația GCTR}
        \label{fig:aes_gctr}
    \end{minipage}\hfill
    \begin{minipage}[t]{0.48\textwidth}
        \centering
        \includegraphics[width=\textwidth]{ghash.png}
        \caption{Generarea etichetei GHASH}
        \label{fig:aes_ghash}
    \end{minipage}
\end{figure}

\vspace{0.5cm}

\paragraph{Funcția GHASH}
Eticheta este calculată prin operații pe corpul finit Galois $\text{GF}(2^{128})$. Fie un șir de blocuri de 128 de biți $X_1, X_2, \dots, X_m$ și subcheia de hash $H = E_K(0^{128})$, funcția se definește iterativ:
$$ Y_i = (Y_{i-1} \oplus X_i) \cdot H $$
unde $Y_0 = 0^{128}$, și rezultatul $Y_m = \text{GHASH}_H(X)$. Înmulțirea se realizează modulo polinomul ireductibil $P(x) = x^{128} + x^7 + x^2 + x + 1$.

\paragraph{Funcția GCTR}
Criptarea este realizată prin modul \textit{CTR}. Pentru fiecare bloc de text în clar $X_i$ și valoarea contorului $CB_i$, textul cifrat $Y_i$ se calculează:
$$ Y_i = X_i \oplus E_K(CB_i) $$
unde $E_K$ este aplicarea algoritmului AES cu secretul $K$.

\paragraph{Criptarea autentificată}
Pentru a proteja atât datele adiționale ($A$), cât și textul cifrat ($C$), eticheta de autentificare finală $T$ (de lungime $t$) este calculată prin aplicarea funcției GHASH asupra concatenării acestora, urmată de o mascare XOR cu valoarea inițială a contorului criptat ($J_0$):
$$ T = \text{MSB}_t \left( \text{GHASH}_H(A \parallel 0^v \parallel C \parallel 0^u \parallel \text{len}(A)_{64} \parallel \text{len}(C)_{64}) \oplus E_K(J_0) \right) $$

\paragraph{Decriptarea autentificată}
La primire, procesul de decriptare verifică autenticitatea datelor înainte de a returna textul în clar, pentru prevenția atacurilor de tip \textit{chosen-ciphertext}. Eticheta anticipată $T'$ este recalculată, utilizând datele adiționale $A$ și textul cifrat primit $C$:
$$ T' = \text{MSB}_t \left( \text{GHASH}_H(A \parallel 0^v \parallel C \parallel 0^u \parallel \text{len}(A)_{64} \parallel \text{len}(C)_{64}) \oplus E_K(J_0) \right) $$

Validarea integrității și extragerea textului în clar $P$ se fac în baza condiției:
$$ P = 
\begin{cases} 
\text{GCTR}_K(J_0, C), & \text{dacă } T = T' \\ 
\mathit{FAIL}, & \text{dacă } T \neq T' 
\end{cases} 
$$

Dacă verificarea eșuează, se întoarce o eroare, iar blocurile decriptate sunt distruse din memorie, garantându-se că pachetele malformate pe rețea sunt respinse imediat.